\section{What a Single Schema Couples}
\label{sec:ch01-what-a-single-schema-couples}

A single GraphQL schema is a single deployable. Every field in it resolves in
one process, so every change to it ships on one release, reviewed by whoever
owns that release. Nothing about GraphQL causes this. It is what happens to any
interface that many teams need to change and one team compiles.

\index{federation!definition}%
Federation breaks the coupling by moving the merge, and the merge here is not
the merge request Ernst's team was reviewing but the act of combining every
team's fields into one schema. Apollo, which publishes the
specification this book builds against, describes it as \enquote{an
architecture for declaratively composing APIs into a unified graph}.
\enquote{Each team can own their slice of the graph independently, empowering
them to deliver autonomously and incrementally}~\autocite{apollofederation}.
\index{subgraph}%
\index{supergraph}%
\index{router}%
Each service publishes its own schema. A composer merges those schemas into one
supergraph. A router serves the supergraph and calls the services to answer
each part of a query.

Netflix says as much about why they moved: \enquote{We still wanted to keep the
unified GraphQL schema of Studio API but decentralize the implementation of the
resolvers to their respective domain teams}~\autocite{shikhare2020}. The schema
the client sees does not change. Who is allowed to edit it does.

That is the whole trade. You give up a merge performed by people and get a
merge performed by a program, and everything federation costs follows from
making that swap.

Nor is this one vendor's idea any longer, which matters if you are deciding
whether to build on it. The GraphQL Specification Working Group announced in
2024 that its Composite Schemas Subcommittee had reconvened and was
\enquote{making steady progress toward a common specification describing
composition and distributed execution across multiple collaborative GraphQL
services}. Engineers from Apollo GraphQL, ChilliCream, Google, Hasura and IBM
were among those contributing to it~\autocite{compositeschemas2024}. This book
builds on Apollo Federation v2 rather than on that draft, and
chapter~\ref{ch:federation-model} says why. What matters here is narrower: the
architecture has outgrown the company that named it, so betting on it is not a
bet on one vendor's roadmap.
